\section{Algoritmet për diagnostikimin e gabimeve numerike}
Qëllimi praktik është të dallohen katër burime: gabimi i të dhënave, kushtëzimi i problemit, gabimi i algoritmit dhe rrumbullakimi. Një diagnozë e mirë ndryshon vetëm një faktor në çdo provë.

\subsection{Matja e saktësisë së makinës}
\begin{lstlisting}[language=Python,caption={Vlerësimi empirik i distancës pranë njësisë.}]
eps = 1.0
while 1.0 + eps/2.0 != 1.0:
    eps /= 2.0
print(eps, np.finfo(float).eps)
\end{lstlisting}
Vlera e marrë është epsilon-i i makinës sipas konventës së NumPy. Njësia e rrumbullakimit është afërsisht gjysma e saj.

\subsection{Shuma Kahan}
Në mbledhjen e shumë termave, një kompensim ruan pjesën e humbur nga rrumbullakimi.
\begin{lstlisting}[language=Python,caption={Mbledhja e kompensuar Kahan.}]
def shuma_kahan(x):
    s = 0.0
    c = 0.0
    for value in x:
        y = value - c
        t = s + y
        c = (t - s) - y
        s = t
    return s

x = np.r_[1e16, np.ones(1_000_000), -1e16]
print(np.sum(x), shuma_kahan(x))
\end{lstlisting}
Ky test duhet përsëritur me renditje të ndryshme. Rezultati tregon se asociativiteti algjebrik nuk ruhet plotësisht në pikë të lëvizshme.

\subsection{Kushtëzimi dhe mbetja}
\begin{lstlisting}[language=Python,caption={Krahasimi i gabimit me mbetjen.}]
def diagnostiko(A, b, x_ref):
    x = np.linalg.solve(A, b)
    mbetja = np.linalg.norm(b - A @ x) / np.linalg.norm(b)
    gabimi = np.linalg.norm(x - x_ref) / np.linalg.norm(x_ref)
    return np.linalg.cond(A), mbetja, gabimi

for e in [1e-2, 1e-6, 1e-10]:
    A = np.array([[1.0, 1.0], [1.0, 1.0 + e]])
    x_ref = np.array([1.0, 1.0])
    b = A @ x_ref
    print(e, diagnostiko(A, b, x_ref))
\end{lstlisting}
Mbetja mund të mbetet e vogël edhe kur gabimi i zgjidhjes rritet. Prandaj raportohen të dyja bashkë me $\kappa(A)$.

\subsection{Algoritmi i auditimit}
\begin{enumerate}
\item Përcaktohet një zgjidhje reference me saktësi të lartë ose analitike.
\item Ndryshohet hapi, toleranca ose rendi i veprimeve.
\item Vizatohet gabimi në shkallë logaritmike.
\item Kontrollohet pjerrësia e pritur e konvergjencës.
\item Përsëritet me të dhëna pak të perturbura për të matur kushtëzimin.
\end{enumerate}

Terminologjia \emph{gabim i përparmë}, \emph{gabim i prapëm}, \emph{kushtëzim} dhe \emph{qëndrueshmëri} përdoret sipas traditës së analizës numerike në shqip, veçanërisht në tekstet e Hanellit--Osmanit dhe Hoxhës.

Për trajtimin në shqip të gabimeve, sistemeve lineare dhe qëndrueshmërisë janë përdorur \cite{hanelli,hoxha,datja}; analiza moderne e qëndrueshmërisë krahasohet me \cite{higham}.
